\documentclass[11pt,a4paper]{article}
\usepackage[T1]{fontenc}
\usepackage[utf8]{inputenc}
\usepackage[margin=25mm]{geometry}
\usepackage{amsmath,amssymb,amsthm,mathtools}
\usepackage{array,booktabs,longtable,enumitem}
\usepackage[hidelinks]{hyperref}
\usepackage{xurl}
\setlength{\emergencystretch}{3em}
\setlength{\parskip}{4pt}
\newtheorem{theorem}{Theorem}[section]
\newtheorem{lemma}[theorem]{Lemma}
\newtheorem{proposition}[theorem]{Proposition}
\newtheorem{corollary}[theorem]{Corollary}
\theoremstyle{definition}
\newtheorem{remark}[theorem]{Remark}
\newcommand{\R}{\mathbb R}
\newcommand{\C}{\mathbb C}
\newcommand{\N}{\mathbb N}
\newcommand{\norm}[1]{\left\lVert#1\right\rVert}
\newcommand{\abs}[1]{\left|#1\right|}
\newcommand{\inner}[2]{\left\langle#1,#2\right\rangle}
\DeclareMathOperator{\diag}{diag}
\DeclareMathOperator{\tr}{tr}
\DeclareMathOperator{\rank}{rank}
\DeclareMathOperator{\supp}{supp}
\DeclareMathOperator{\Span}{span}
\newcommand{\code}[1]{\ifmmode\text{\small\nolinkurl{#1}}\else{\small\nolinkurl{#1}}\fi}
\title{Four focused mathematical reroutes\\
Explicit reciprocal denominators, finite degree transfer,\\
Jacobi borders, and rational measure modification}
\author{Mathematical handoff for Lean formalization}
\date{14 October 2026}
\begin{document}
\maketitle

This standalone note contains only the four requested targets. It gives
mathematical derivations and identifies reusable statements; it does not
implement or claim new Lean proofs.

\begin{enumerate}[leftmargin=*]
\item \textbf{Reciprocal exponential approximation.} An explicit degree-$n$
denominator is positive and strictly increasing on the half-line for $n\ge1$.
Its uniform error is at most $\sqrt{40/21}\,3^{-n}$, which retains the earlier
$\sqrt2\,3^{-n}$ inequality. Its construction needs no weighted minimizer.
The lower bound needed for the exact root limit $1/3$ is retained and proved
separately in Appendix~\ref{app:retained-lower}; that proof still uses weighted
minimax and alternation.
\item \textbf{The polynomial Bernstein limit for $|x|$.} A finite inequality
transfers a sufficiently large source degree to any target degree. A scalar
liminf then proves convergence. No limiting entire function or coefficient
subsequence is required. Finite complex polynomial growth on ellipses remains
an explicit, already available input.
\item \textbf{Radau and Lobatto borders.} A telescoping matrix-power identity
proves locality and the exact moment-preservation thresholds. Scalar Schur
complements determine the borders and keep their spectra in the interval.
\item \textbf{Christoffel--Uvarov modification.} A real coefficient system
proves existence and the nonzero normalization determinant for every
denominator degree, including $m>n$. Polynomial remainders also cover repeated
roots. The false complex-point nonvanishing conclusion is explicitly corrected.
\end{enumerate}

All sums defining the new denominators and coefficient systems are finite.
Empty matrices have determinant one. Every normalization or inverse below
comes with its required nonvanishing or positivity hypothesis. The proof
interfaces are distinguished from the mathematics they express.

The public declarations cited below were checked in the saved index at
\href{https://github.com/tchen-research/NLAlib/tree/e5725bb94e5bd8e326d58ad219260b0218e74c9e}
{commit \texttt{e5725bb94e5b}}. The current local index is different and omits
some of those declarations; reuse is therefore qualified by that snapshot.
No build or new kernel verification is claimed.

\tableofcontents
\clearpage


\section{Reciprocal exponential approximation: an explicit increasing denominator}
\label{sec:explicit-reciprocal}

The upper-bound construction can be made completely explicit. Two finite
polynomial resolvents replace the weighted minimizer, alternation, and
extrapolation used to construct the previous denominator. The new bound is
$\sqrt{40/21}\,3^{-n}\le\sqrt2\,3^{-n}$, so it retains the previously
stated constant as well as the sharp root rate.
The previously proved lower bound is still needed to identify the root limit.
We do not claim that the construction below replaces that lower-bound proof.

\subsection{The denominator and its quantitative statement}

Let $\mathcal P_n$ be the real polynomials of degree at most $n$, and let
$D:\mathcal P_n\to\mathcal P_n$ be differentiation. Define
\[
 L_j(x)=\sum_{r=0}^j(-1)^r\binom jr\frac{x^r}{r!},\qquad
 g_n(x)=\sum_{j=0}^n\frac{4(-1)^j}{3^{j+1}}L_j(x).
\]
For $n\ge1$ put
\begin{equation}\label{eq:explicit-denominator}
 H_n(x)=\frac{15}{26}\sum_{r=0}^n
       \left((4/3)^{r+1}-4^{-r-1}\right)D^r g_n(x),
 \qquad Q_n(t)=H_n(4t),
\end{equation}
and set $Q_0(t)=1$. Thus all coefficients of $Q_n$ are explicit rational
numbers obtained from finite sums. Some monomial coefficients can be negative;
coefficientwise positivity is neither asserted nor needed.

\begin{theorem}\label{thm:positive-reciprocal}
For $n\ge1$, $Q_n$ has degree exactly $n$ and satisfies
\[
 Q_n(t)>0,\qquad Q_n'(t)>0\qquad(t\ge0).
\]
For every integer $n\ge0$,
\begin{equation}\label{eq:explicit-reciprocal-error}
 \sup_{t\ge0}\left|e^{-t}-\frac1{Q_n(t)}\right|
 \le\sqrt{\frac{40}{21}}\,3^{-n}\le\sqrt2\,3^{-n}.
\end{equation}
\end{theorem}

The following proof needs finite polynomial identities, exponential integrals,
the spectral theorem for a real symmetric tridiagonal matrix, and scalar
Cauchy--Schwarz. It uses no minimizing polynomial.

\subsection{Finite orthogonality and the exact squared error}

Write $f(x)=e^{x/4}$ and
\[
 \langle u,v\rangle=\int_0^\infty u(x)v(x)e^{-x}\,dx.
\]
The finite Rodrigues identity
\[
 L_j(x)=\frac{e^x}{j!}D^j(e^{-x}x^j)
\]
and $j$ integrations by parts give
\[
 \langle L_j,p\rangle
 =\frac{(-1)^j}{j!}\int_0^\infty x^j e^{-x}D^jp(x)\,dx.
\]
The boundary terms vanish: derivatives of $e^{-x}x^j$ of order below $j$
vanish at zero, and exponential decay handles infinity. For a polynomial $p$
of degree less than $j$ the result is zero. Since the leading coefficient of
$L_j$ is $(-1)^j/j!$, the identity also gives $\langle L_j,L_j\rangle=1$.
The same integrations by parts with $p=f$, using
$\int_0^\infty x^j e^{-s x}\,dx=j!/s^{j+1}$, give
\[
 \langle f,L_j\rangle=\frac{4(-1)^j}{3^{j+1}},
 \qquad \langle f,f\rangle=2.
\]
Expanding a square and summing a finite geometric progression now proves
\begin{equation}\label{eq:laguerre-finite-error}
 \int_0^\infty |f(x)-g_n(x)|^2e^{-x}\,dx
 =2-\sum_{j=0}^n\frac{16}{9^{j+1}}
 =\frac2{3^{2n+2}}.
\end{equation}
No infinite orthogonal expansion or completeness theorem is involved.

\subsection{Two smoothing steps control both value and derivative}

For $a>0$, nilpotence $D^{n+1}=0$ gives the finite identity
\[
 R_a=(aI-D)^{-1}=\sum_{r=0}^n a^{-r-1}D^r.
\]
For a polynomial $p$, repeated integration by parts yields
\[
 (R_ap)(x)=\int_0^\infty e^{-au}p(x+u)\,du.
\]
It can also be verified termwise from the finite Taylor formula and
$\int_0^\infty u^re^{-au}\,du=r!/a^{r+1}$. The elementary resolvent identity
is $(b-a)R_aR_b=R_a-R_b$. Therefore
\[
 H_n=\frac{15}{8}R_{3/4}R_4g_n
     =\frac{15}{26}(R_{3/4}-R_4)g_n.
\]
Put
\[
 k(u)=\frac{15}{26}(e^{-3u/4}-e^{-4u}),\qquad
 \ell(u)=-k'(u)
       =\frac{15}{26}\left(\frac34e^{-3u/4}-4e^{-4u}\right).
\]
Then
\[
 H_n(x)=\int_0^\infty k(u)g_n(x+u)\,du,\qquad
 H_n'(x)=\int_0^\infty\ell(u)g_n(x+u)\,du.
\]
The second identity is integration by parts, using $k(0)=0$. These formulas
require no differentiation under an improper integral: $H_n$ is already
the finite polynomial in \eqref{eq:explicit-denominator}. Direct integration
also gives
\[
 \int_0^\infty k(u)e^{u/4}\,du=1,\qquad
 \int_0^\infty\ell(u)e^{u/4}\,du=\frac14.
\]
Thus the two kernels applied to $f(x+u)$ give $f(x)$ and $f(x)/4$.

For $x\ge0$, the weighted squared error on $[x,\infty)$ is bounded by
\eqref{eq:laguerre-finite-error}. Apply Cauchy--Schwarz to each kernel,
using
\[
 \int_0^\infty k(u)^2e^u\,du=\frac{15}{28},\qquad
 \int_0^\infty\ell(u)^2e^u\,du=\frac{135}{224}.
\]
For example, the first integral equals
$\frac{225}{676}(2-8/15+1/7)=15/28$; the second equals
$\frac{225}{676}(9/8-8/5+16/7)=135/224$.
This gives the two simultaneous estimates
\begin{align}
 |H_n(x)-f(x)|
 &\le C\,3^{-n}e^{x/2},\qquad C=\sqrt{\frac5{42}},
 \label{eq:twice-smooth-value}\\
 |H_n'(x)-f(x)/4|
 &\le C'\,3^{-n}e^{x/2},\qquad C'=\sqrt{\frac{15}{112}}.
 \label{eq:twice-smooth-derivative}
\end{align}
For clarity, the factor $e^{x/2}$ comes from
\[
 \int_0^\infty |f(x+u)-g_n(x+u)|^2e^{-u}\,du
 =e^x\int_x^\infty |f(t)-g_n(t)|^2e^{-t}\,dt.
\]

\subsection{A finite tridiagonal bound proves positivity beyond the cutoff}

First consider $0\le x\le4n$, with $n\ge5$. Dividing
\eqref{eq:twice-smooth-value}--\eqref{eq:twice-smooth-derivative} by $f(x)$
gives
\[
 H_n(x)\ge f(x)\left(1-C(e/3)^n\right)>0,
\]
\[
 H_n'(x)\ge\frac{f(x)}4
   \left(1-4C'(e/3)^n\right)>0.
\]
For the value, use $C<1$ and $e/3<1$. For the derivative, square its
nonnegative error factor and use
\[
 16(C')^2(e/3)^{2n}
 <\frac{15}{7}(11/12)^{10}
 =\frac{129687123005}{144473849856}<1.
\]
Here $e<11/4$ follows from its power series: for $k\ge3$,
$k!\ge6\,3^{k-3}$, with strict inequality for $k\ge4$, so the tail is
strictly less than $1/4$ after the first three terms $1+1+1/2$.

The four remaining positive degrees are checked by finite rational
arithmetic. In the table, $H_n(x)=d_n^{-1}\sum_{r=0}^n a_rx^r$:
\[
\begin{array}{c|r|rrrrr}
 n&d_n&a_0&a_1&a_2&a_3&a_4\\\hline
 1&216&215&60&&&\\
 2&7776&7745&1860&360&&\\
 3&279936&280055&70620&6840&1440&\\
 4&10077696&10075985&2513220&341640&10080&4320
\end{array}
\]
Every displayed coefficient is positive; hence both $H_n$ and $H_n'$
are strictly positive on $[0,\infty)$ for $1\le n\le4$.

For $x\ge4n$, use the location of the roots of the finitely many $L_j$.
Here a deliberately coarse bound is enough. Put
$P_j=(-1)^j j!L_j$. Direct comparison of coefficients gives
\[
 P_0=1,\quad P_1=x-1,\quad
 P_{j+1}=(x-(2j+1))P_j-j^2P_{j-1}.
\]
Consequently $P_j(x)=\det(xI-J_j)$, where $J_j$ is the $j\times j$
real symmetric matrix indexed by $r=0,\ldots,j-1$ with
\[
 (J_j)_{rr}=2r+1,\qquad
 (J_j)_{r,r+1}=(J_j)_{r+1,r}=r+1.
\]
This determinant recurrence is just expansion along the last row.
For any real vector $v$, the inequality $2ab\le a^2+b^2$ gives
\[
 v^\top J_jv\le(4j-2)\|v\|_2^2\qquad(j\ge1).
\]
Indeed, the coefficient of $v_r^2$ after this estimate is at most
$(2r+1)+r+(r+1)=4r+2\le4j-2$; missing boundary terms only decrease it.
Thus all the roots $\lambda_{j,1},\ldots,\lambda_{j,j}$ of $P_j$ are
real and satisfy $\lambda_{j,k}\le4j-2<4j$. For $1\le j\le n$ and
$x\ge4n$,
\[
 \frac{4(-1)^j}{3^{j+1}}L_j(x)
 =\frac4{3^{j+1}j!}\prod_{k=1}^j(x-\lambda_{j,k}).
\]
Every factor is strictly positive. The product rule shows that every
derivative through order $j$ is strictly positive there, and higher
derivatives vanish. The $j=0$ term is the positive constant $4/3$.
It follows that $g_n(x)>0$, that all its derivatives are nonnegative,
and that $g_n'(x)>0$ when $n\ge1$.
Every coefficient multiplying $D^rg_n$ in
\eqref{eq:explicit-denominator} is positive. That finite expression proves
$H_n(x)>0$ and $H_n'(x)>0$ for $x\ge4n$.

This completes the positivity proof on the whole half-line. Notice that
the cutoff $4n$ is used only to prove monotonicity. The approximation-error
cutoff below is different. Finally the coefficient of $x^n$ in $H_n$ is
$\frac58\frac4{3^{n+1}n!}>0$, proving its degree assertion.

\subsection{Uniform reciprocal error and the retained lower bound}

For $n\ge1$ set
\[
 \delta_n=C\,3^{-n},\qquad
 a_n=4n\log3+4\log\frac1{2C}>0.
\]
Then $\delta_ne^{a_n/4}=1/2$. On $[0,a_n]$,
\eqref{eq:twice-smooth-value} implies $H_n\ge f/2$, so
\[
 \left|f(x)^{-1}-H_n(x)^{-1}\right|
 =\frac{|H_n(x)-f(x)|}{f(x)H_n(x)}
 \le2\delta_n.
\]
For $x\ge a_n$, monotonicity and the bound at $a_n$ give
\[
 H_n(x)^{-1}\le H_n(a_n)^{-1}\le2e^{-a_n/4}
 =4\delta_n,\qquad
 f(x)^{-1}\le e^{-a_n/4}=2\delta_n.
\]
Both quantities are nonnegative; hence their difference in absolute value
is at most their maximum, proving \eqref{eq:explicit-reciprocal-error}
after $x=4t$, since $4C=\sqrt{40/21}<\sqrt2$. For $n=0$, $Q_0=1$
has error $1\le\sqrt{40/21}$ directly.

Define $E_{0,n}$ as the infimum of the uniform errors of real rational
functions of type $(0,n)$ with no pole on $[0,\infty)$. The lower estimate
from the earlier document, reproduced in Appendix~\ref{app:retained-lower}, is
\[
 E_{0,n}(e^{-t};[0,\infty))\ge
 \frac{3^{-n}}{6\sqrt{(4n+4)\log3+\log4}}.
\]
Together with the explicit upper bound above it gives
\[
 \lim_{n\to\infty}E_{0,n}(e^{-t};[0,\infty))^{1/n}=\frac13.
\]
\emph{Dependency retained:} the cited lower estimate uses the previous
weighted-minimax comparison and the finite Gram identity. The Gram identity
by itself is an $L^2$ polynomial-approximation statement and does not directly
give a lower bound for uniform reciprocal approximation. The present route
fully replaces the upper-bound construction, but does not silently claim to
remove that separate lower-bound argument.

\paragraph{Operator form and formalization interface.}
For every finite-dimensional real symmetric or complex Hermitian positive
semidefinite operator $A$, diagonalization gives
\[
 Q_n(A)\succ0,\qquad
 \|e^{-A}-Q_n(A)^{-1}\|_2
 \le\sqrt{\frac{40}{21}}\,3^{-n}\le\sqrt2\,3^{-n}.
\]
The denominator definition is a finite polynomial sum. The new reusable
lemmas are the two polynomial-resolvent integral identities, their two
scalar error estimates, and a coarse spectral bound for $J_j$.
The declaration index was searched for Laguerre and reciprocal-exponential
approximation results; it does not currently supply this construction.
The tridiagonal determinant and spectral steps can reuse the existing
Jacobi-matrix and symmetric-matrix interfaces. No weighted minimizer,
alternation theorem, or extrapolation theorem is required for this upper
bound.


\section{The Bernstein limit for \texorpdfstring{$|x|$}{abs(x)} by a finite degree-transfer inequality}
\label{sec:finite-bernstein}

\paragraph{Outcome.}
There is a route that removes the limiting entire function altogether.
For a target degree $N$, restrict and rescale a best polynomial of a much
larger degree $n$, then discard its Chebyshev coefficients above $N$.
A bound uniform in the source degree $n$ gives a comparison of the two
finite errors. Taking a scalar liminf proves convergence. No normal-family
theorem, diagonal extraction of coefficients, or compactness of a space
of entire functions is used. The one existing complex-analytic input is
the finite polynomial growth estimate on a Bernstein ellipse; even the
new change-of-variables estimate below uses only real inverse functions.

Write
\[
 E_n=\inf_{\deg p\le n}\max_{-1\le x\le1}|\,|x|-p(x)|.
\]
All polynomials in this definition are real. Their complex evaluations
will be used only for finite inequalities.

\subsection{Two finite polynomial facts}
For $\rho>1$, write
\[
 \mathcal E_\rho=
 \left\{(w+w^{-1})/2:|w|=\rho\right\}
\]
for the Bernstein ellipse, and let its filled version include the interior.
We use the already available polynomial growth inequality
\begin{equation}
 \deg p\le n,\quad \|p\|_{[-1,1]}\le M
 \quad\Longrightarrow\quad
 |p(z)|\le M\rho^n\quad
 (z\text{ in the filled ellipse}).
 \label{eq:bern-finite-growth}
\end{equation}
At $\rho=1$, the assertion simply means the original interval bound.
One proof applies the maximum-modulus principle to the ordinary
polynomial $v^np((v+v^{-1})/2)$ on $|v|\le1$, and substitutes $v=w^{-1}$.
There is no limiting family of analytic functions in this fact.

The companion truncation statement can be proved with only finite Fourier
orthogonality. If a real polynomial $r$ has $|r(z)|\le M$ on
$\mathcal E_\rho$, then there is a real polynomial $q$ of degree at most $N$
such that
\begin{equation}
 \|r-q\|_{[-1,1]}\le
 \frac{2M\rho^{-N}}{\rho-1}.
 \label{eq:bern-finite-truncation}
\end{equation}
Indeed, the finite Laurent expansion is
\[
 r((w+w^{-1})/2)=\sum_{k=-d}^{d}c_kw^k,
 \qquad c_{-k}=c_k\in\R.
\]
Integrating this finite identity against $e^{-ik\theta}$ on
$w=\rho e^{i\theta}$ gives $|c_k|\le M\rho^{-k}$ for $k\ge0$.
Take
$q(x)=c_0+2\sum_{k=1}^{\min(N,d)}c_kT_k(x)$, use
$|T_k(x)|\le1$ on $[-1,1]$, and sum a finite geometric tail.
This proves \eqref{eq:bern-finite-truncation} without an infinite
Chebyshev series or a theorem about analytic approximation.

\subsection{A small-coordinate lemma using real arithmetic}
\begin{lemma}\label{lem:bern-sine-coordinates}
If $\zeta=x+iy$ and $r=|\zeta|\le1/2$, there exist real $a,b$ such that
\[
 \sin(a+ib)=\zeta,\qquad
 |a+ib|\le2r,\qquad
 |b|\le\frac{|y|}{\sqrt{1-r^2}}.
\]
Consequently $\zeta$ belongs to the filled Bernstein ellipse of parameter
$e^{|b|}$, with the interval understood when $b=0$.
\end{lemma}
\begin{proof}
If $y=0$, take $b=0$ and $a=\arcsin x$ using the real inverse sine.
If $y\ne0$, put
\[
 v=\frac{\sqrt{(1-r^2)^2+4y^2}-(1-r^2)}2>0.
\]
Then
\[
 v(v+1-r^2)=y^2,
 \qquad \frac{x^2}{1+v}+\frac{y^2}{v}=1,
 \qquad v\le\frac{y^2}{1-r^2}.
\]
Set
\[
 a=\arcsin\!\left(\frac{x}{\sqrt{1+v}}\right),
 \qquad b=\operatorname{sgn}(y)\operatorname{arsinh}\sqrt v.
\]
The real formulas for sine, cosine, hyperbolic sine and hyperbolic cosine
give $\sin a\cosh b=x$ and $\cos a\sinh b=y$. Thus
$\sin(a+ib)=x+iy$. The real derivative bounds
$|\arcsin u|\le2|u|$ for $|u|\le1/2$ and
$\operatorname{arsinh}u\le u$ for $u\ge0$ give
\[
 |a|\le2|x|,\qquad
 |b|\le\sqrt v\le\frac{|y|}{\sqrt{1-r^2}}
 \le2|y|.
\]
Their squares imply $|a+ib|\le2r$.
Finally $\sin(a+ib)=(w+w^{-1})/2$ with
$w=i e^{-i(a+ib)}$; either $w$ or $w^{-1}$ has modulus $e^{|b|}$.
\end{proof}

This lemma deliberately avoids a complex inverse-sine function, its
branch cut, and a theorem about convergence of its rescalings.

\subsection{The uniform finite bound after subtracting a comparison polynomial}
For even $n=2m\ge2$, use the explicit polynomial
\[
 S_n(x)=\frac2\pi+\frac4\pi\sum_{j=1}^{m}
      \frac{(-1)^{j+1}T_{2j}(x)}{4j^2-1}.
\]
The known absolutely convergent Chebyshev expansion of $|x|$ gives
\begin{equation}
 \|S_n-|x|\|_{[-1,1]}=\frac2{\pi(n+1)},
 \qquad S_n(0)=\frac2{\pi(n+1)}.
 \label{eq:bern-comparison}
\end{equation}
The norm upper bound is the telescoping coefficient tail, and equality
holds at zero since the omitted terms there have the same sign.

Choose a minimizer $p_n$, whose existence is a finite-dimensional
compactness result already available in the library, and put
\[
 P_n(z)=np_n(z/n),\quad B_n(z)=nS_n(z/n),\quad D_n=P_n-B_n.
\]
Since $E_n\le2/(\pi(n+1))$, we have
\begin{equation}
 \|D_n\|_{[-n,n]}\le nE_n+\frac{2n}{\pi(n+1)}
 \le\frac4\pi.
 \label{eq:bern-bounded-difference}
\end{equation}

\begin{lemma}[Uniform bound with an explicit source-degree threshold]
\label{lem:bern-uniform-finite}
Let $R\ge0$ and let $n\ge2$ be even, with
\[
 n\ge2\max(1,R^2).
\]
For every $z\in\C$ with $|z|\le R$,
\begin{equation}
 |P_n(z)|\le\frac{6(1+R^2)}\pi
                   e^{|\operatorname{Im}z|+1}.
 \label{eq:bern-uniform-finite}
\end{equation}
\end{lemma}
\begin{proof}
The threshold implies $R/n\le1/2$. Apply
Lemma~\ref{lem:bern-sine-coordinates} to $z/n$, obtaining $u=a+ib$.
Equation~\eqref{eq:bern-finite-growth}, applied to
$D_n(n\,\cdot)$, gives
\[
 |D_n(z)|\le(4/\pi)e^{n|b|}.
\]
The finite Chebyshev recurrence gives
$T_{2j}(\sin u)=(-1)^j\cos(2ju)$. Therefore
\[
 B_n(z)=\frac{2n}{\pi(n+1)}+
       \frac{4n}{\pi}\sum_{j=1}^{n/2}
       \frac{1-\cos(2ju)}{4j^2-1}.
\]
Twice integrating the derivative of cosine over a real parameter gives
\[
 1-\cos w=w^2\int_0^1(1-s)\cos(sw)\,ds,
 \qquad |1-\cos w|\le\frac{|w|^2}{2}e^{|\operatorname{Im}w|}.
\]
Since $j^2/(4j^2-1)\le1/3$ for $j\ge1$ and $|u|\le2|z|/n$,
\[
 |B_n(z)|\le\frac2\pi+
     \frac{4n^2}{3\pi}|u|^2e^{n|b|}
 \le\frac{2+(16/3)R^2}{\pi}e^{n|b|}.
\]
Adding the bound for $D_n$ yields
$|P_n(z)|\le6(1+R^2)e^{n|b|}/\pi$.

For $0\le s\le1/4$, $(1-s)^{-1/2}\le1+2s$.
Lemma~\ref{lem:bern-sine-coordinates} and the threshold now imply
\[
 n|b|\le\frac{|\operatorname{Im}z|}{\sqrt{1-R^2/n^2}}
 \le |\operatorname{Im}z|+\frac{2R^3}{n^2}
 \le |\operatorname{Im}z|+1.
\]
The final inequality follows separately for $R\le1$ and $R\ge1$;
in fact the last added term is at most $1/2$.
\end{proof}

\subsection{A finite comparison of two degrees}
\begin{theorem}[Degree transfer]\label{thm:bern-degree-transfer}
Fix $0<\alpha<1$ and $N\ge1$. Set
\[
 t=\frac{1-\alpha}{2},\quad
 R=\alpha N\cosh t,\quad c=t(1-e^{-t})>0.
\]
For every even $n\ge2$ satisfying $n\ge2\max(1,R^2)$,
\begin{equation}
 NE_N\le\frac{nE_n}{\alpha}+
 \frac{12e(1+\alpha^2N^2\cosh^2t)}
      {\pi\alpha(e^t-1)}e^{-cN}.
 \label{eq:bern-degree-transfer}
\end{equation}
\end{theorem}
\begin{proof}
Apply \eqref{eq:bern-finite-truncation}, with $\rho=e^t$, to the
finite polynomial
\[
 r(x)=\frac{P_n(\alpha Nx)}{\alpha N}
     =\frac{n}{\alpha N}p_n\!\left(\frac{\alpha Nx}{n}\right).
\]
On the ellipse, $|\alpha Nz|\le R$ and
$|\operatorname{Im}(\alpha Nz)|\le\alpha N\sinh t$.
Lemma~\ref{lem:bern-uniform-finite} gives the bound
\[
 M=\frac{6(1+R^2)}{\pi\alpha N}
           e^{\alpha N\sinh t+1}.
\]
On the real interval the argument of $p_n$ lies in $[-1,1]$, since
$\alpha N\le R\le n/2$. Hence
$\|r-|x|\|\le nE_n/(\alpha N)$. Adding the truncation error and
multiplying by $N$ gives the asserted inequality with exponent
$N(\alpha\sinh t-t)$. Finally
\[
 \alpha=1-2t\le e^{-2t},\qquad
 \sinh t\le te^t,\qquad
 \alpha\sinh t-t\le-t(1-e^{-t})=-c.
\]
\end{proof}

\subsection{Convergence uses only a scalar liminf}
Let $b=\liminf_{m\to\infty}(2m)E_{2m}$.
For every fixed $N$, every $0<\alpha<1$, and every $\eta>0$, there are
arbitrarily large even $n$ with $nE_n\le b+\eta$.
Choose one above the explicit threshold in
Theorem~\ref{thm:bern-degree-transfer}. Its error term does not depend on
$n$. Letting $\eta\downarrow0$ therefore gives
\[
 NE_N\le\frac b\alpha+
 \frac{12e(1+\alpha^2N^2\cosh^2t)}
      {\pi\alpha(e^t-1)}e^{-cN}.
\]
The second term tends to zero as $N\to\infty$, so
$\limsup NE_N\le b/\alpha$. Let $\alpha\uparrow1$.
Since the full limsup dominates the even liminf $b$, this already forces
the even subsequence to converge to $b$. Symmetrization gives
$E_{2m+1}=E_{2m}$, and
$(2m+1)E_{2m+1}=(1+1/(2m))(2m)E_{2m}$, so the full sequence converges:
\begin{equation}
 \lim_{n\to\infty}nE_n=\beta=b.
 \label{eq:bern-limit}
\end{equation}

For completeness, positivity has a finite dual certificate. Put
$K=\lfloor n/2\rfloor+1$, $w_j=1-|j|/(K+1)$, and
\[
 F_K(t)=\frac1{K+1}\left|\sum_{j=0}^{K}e^{ijt}\right|^2,
 \qquad
 \Lambda(h)=\frac1{2\pi}\int_{-\pi}^{\pi}
       h(\sin t)\cos(4Kt)F_K(2t)\,dt.
\]
The kernel $F_K$ is nonnegative with normalized integral one, hence
$\|\Lambda\|\le1$. Every frequency of its modulated version has
absolute value at least $2K>n$, so $\Lambda(p)=0$ for $\deg p\le n$.
Finite Fourier orthogonality and direct integration of $|\sin t|$ give
\[
 E_n\ge|\Lambda(|x|)|
 =\sum_{j=-K}^{K}\frac{2w_j}{\pi(4(2K+j)^2-1)}
 \ge\frac{2(K+1)}{\pi(36K^2-1)}
 \ge\frac1{18\pi(n+1)}.
\]
Together with \eqref{eq:bern-comparison}, this proves
\[
 \frac1{18\pi}\le\beta\le\frac2\pi.
\]
These coarse enclosing constants establish positivity and finiteness;
no numerical value of the Bernstein constant is asserted here.

\paragraph{Formalization boundary and checked reuse.}
The new centerpiece is the explicit finite inequality
\eqref{eq:bern-degree-transfer}. It replaces all of the comparison-function
limit, local uniform convergence, coefficient subsequence, and entire
function truncation steps of the previous derivation. Besides real
arithmetic, finite trigonometric sums and real-variable integration, it
uses the already proved ellipse growth inequality
\eqref{eq:bern-finite-growth}. Thus it does not pretend to remove complex
evaluation from the proof, but it needs no new complex compactness theory.

The following names were checked against the saved public declaration
index at commit \texttt{e5725bb94e5b}; the current checkout's smaller
index does not list all of them, so their local availability should be
checked when choosing the formalization branch:
\begin{itemize}
\item \code{NLAlib.exists_degree_le_norm_sub_eq_minimaxError} supplies each
      finite-dimensional minimizer.
\item \code{NLAlib.hasSum_chebyshevCoeff_abs_mul_eval_T},
      \code{NLAlib.chebyshevCoeff_abs_even}, and
      \code{NLAlib.tsum_abs_chebyshevCoeff_abs_tail} supply
      \eqref{eq:bern-comparison}.
\item \code{NLAlib.norm_eval_le_mul_pow_of_mem_filledBernsteinEllipse}
      supplies \eqref{eq:bern-finite-growth}.
\item \code{NLAlib.exists_degree_le_abs_sub_eval_le_of_open_bernsteinEllipse}
      can supply \eqref{eq:bern-finite-truncation}; the finite Laurent
      proof above is an alternative requiring a smaller interface.
\item Mathlib's real \code{Real.arcsin}, \code{Real.hasDerivAt_arcsin},
      \code{Real.arsinh}, and \code{Real.sinh_arsinh} are present in the
      checked local Mathlib source. They suffice for the coordinate lemma;
      a complex inverse-sine definition is unnecessary.
\end{itemize}



\section{Radau and Lobatto borders from finite propagation}
\label{focused:borders}

The main simplification is to construct the borders by two scalar Schur
complements and to prove moment preservation by a telescoping identity.
No counting of walks and no limiting argument is needed.
All matrix indices in this section start at zero.

\subsection{A single locality lemma}

\begin{lemma}[Rooted matrix-power locality]
Let $A,B$ be tridiagonal $N\times N$ matrices over a field, with $N\ge1$.
Suppose $A_{ij}=B_{ij}$ whenever $i+j<s$, where $s\ge0$ is an integer.
Then
\[
 (A^k)_{00}=(B^k)_{00}\qquad(0\le k\le s).
 \tag{BL1}
\]
\end{lemma}
\begin{proof}
Induction in $r$ gives
$(A^r)_{0i}=0$ for $i>r$ and $(B^r)_{j0}=0$ for $j>r$.
For $k\ge1$, use the noncommutative telescoping identity
\[
 A^k-B^k=\sum_{r=0}^{k-1}A^r(A-B)B^{k-1-r}.
\]
In the $(0,0)$ entry, a potentially nonzero summand has
$i\le r$, $j\le k-1-r$, and therefore $i+j\le k-1<s$.
The middle factor $(A-B)_{ij}$ is zero. The case $k=0$ is immediate.
\end{proof}

In particular, changing only the diagonal entry at index $q$ preserves
moments through degree $2q$. Adding a coupling between indices $q$ and
$q+1$ preserves them through degree $2q+1$. These thresholds are sharp:
if the earlier off-diagonals are $a_1,\ldots,a_q$, a change $\delta$ in
that diagonal first contributes
$\delta\prod_{j=1}^q a_j^2$ at degree $2q+1$; adding a previously zero
coupling $b$ first contributes $b^2\prod_{j=1}^q a_j^2$ at degree $2q+2$.
Empty products equal one. The same induction on support verifies these
first nonzero contributions.

\subsection{Why a Jacobi compression already has the required moments}

Let $\mu$ be a nonzero finite positive measure supported in $[a,c]$,
$a<c$, and put $h_0=\int1\,d\mu$. Suppose its polynomial inner product
is positive definite through degree $N-1$. Let
$p_0,\ldots,p_{N-1}$ be its orthonormal polynomials, with positive
leading coefficients and $p_0=h_0^{-1/2}$, and define
\[
 (J_N)_{ij}=\int t p_i(t)p_j(t)\,d\mu(t).
\]
Polynomial degrees show that $J_N$ is symmetric tridiagonal, and
\[
 aI\preceq J_N\preceq cI.
 \tag{BL2}
\]
For $0\le r<N$, the coefficient vector of $t^rp_0$ in this basis is
$J_N^re_0$: until this degree, multiplication followed by orthogonal
projection has discarded nothing. Consequently, for $r,s<N$,
\[
 h_0^{-1}\int t^{r+s}\,d\mu
   =(J_N^re_0)^\top(J_N^se_0).
\]
This gives moments through $2N-2$. The one remaining moment is
\[
 h_0^{-1}\int t^{2N-1}\,d\mu
 =\int t(t^{N-1}p_0)^2\,d\mu
 =(J_N^{N-1}e_0)^\top J_N(J_N^{N-1}e_0).
\]
Thus
\[
 \int t^k\,d\mu=h_0(J_N^k)_{00}\qquad(0\le k\le2N-1).
 \tag{BL3}
\]
Only the $N$ displayed orthonormal polynomials are required. In
particular, this proof does not ask for a normalized polynomial of
degree $N$, which might fail at finite-support termination.

\subsection{Radau: one last diagonal and one Schur complement}

Assume $q\ge1$ and the polynomial inner product is positive definite
through degree $q$. Write
\[
 J_{q+1}=\begin{pmatrix}J_q&b e\\b e^\top&\alpha_q\end{pmatrix},
 \qquad e=e_{q-1},\quad b>0.
\]
Here $b=a_q=\sqrt{\beta_q}$ in the monic recurrence convention.
The strict inequalities $aI\prec J_q\prec cI$ follow directly from
the quadratic-form formula: a nonzero polynomial of degree at most
$q-1$ cannot vanish at all the support points other than one endpoint.
Positive definiteness through degree $q$ supplies at least $q+1$
support points.

Define
\begin{align*}
 d_a&=a+b^2 e^\top(J_q-aI)^{-1}e,&
 R_a&=\begin{pmatrix}J_q&b e\\b e^\top&d_a\end{pmatrix},\\
 d_c&=c-b^2 e^\top(cI-J_q)^{-1}e,&
 R_c&=\begin{pmatrix}J_q&b e\\b e^\top&d_c\end{pmatrix}.
\end{align*}
The Schur complement of $J_q-aI$ in $R_a-aI$ is zero. Hence
$R_a-aI$ is positive semidefinite and singular. The same Schur
complement in $J_{q+1}-aI$ is nonnegative, so $d_a\le\alpha_q$.
Similarly $d_c\ge\alpha_q$ and $cI-R_c$ is positive semidefinite
and singular. Combining this with (BL2) yields
\[
 aI\preceq R_a\preceq J_{q+1}\preceq R_c\preceq cI,
 \qquad a\in\operatorname{spec}(R_a),\quad
 c\in\operatorname{spec}(R_c).
 \tag{BL4}
\]
Both borders differ from $J_{q+1}$ only at diagonal index $q$.
Equations (BL1) and (BL3) therefore give
\[
 h_0(R_a^k)_{00}=\int t^k\,d\mu=h_0(R_c^k)_{00}
 \qquad(0\le k\le2q).
 \tag{BL5}
\]

For comparison with the usual formula, set
$\pi_j(t)=\det(tI-J_j)$, $\pi_0=1$.
A cofactor gives
$e^\top(xI-J_q)^{-1}e=\pi_{q-1}(x)/\pi_q(x)$, hence
\[
 d_z=z-\beta_q\frac{\pi_{q-1}(z)}{\pi_q(z)},\qquad z=a,c.
 \tag{BL6}
\]
Expansion in the last row gives
$\det(tI-R_z)=(t-d_z)\pi_q(t)-\beta_q\pi_{q-1}(t)$.
Thus the standard polynomial border and the Schur-complement border
are exactly the same construction.
For $q=0$, use the one-by-one matrices $[a]$ and $[c]$; only the
constant moment is claimed.

\subsection{Lobatto: two endpoint equations}

Let $q\ge0$, set $N=q+1$, and assume that $J_N$ is defined and
\[
 aI\prec J_N\prec cI.
 \tag{BL7}
\]
These are the exact endpoint-invertibility and positivity hypotheses
used below. A clean sufficient measure hypothesis is at least $q+2$
support points in $[a,c]$. Merely $q+1$ support points need not suffice
when the support includes an endpoint.
Put $e=e_q$ and
\[
 r(x)=e^\top(xI-J_N)^{-1}e,
 \qquad
 \gamma=\frac{c-a}{r(c)-r(a)},\qquad
 d=a-\gamma r(a)=c-\gamma r(c).
\]
Since $r(a)<0<r(c)$, one has $\gamma>0$. Define
\[
 B=\begin{pmatrix}J_N&\sqrt\gamma e\\
                   \sqrt\gamma e^\top&d\end{pmatrix}.
 \tag{BL8}
\]
The endpoint Schur complements are both zero:
\[
 d-a-\gamma e^\top(J_N-aI)^{-1}e=0,
 \qquad
 c-d-\gamma e^\top(cI-J_N)^{-1}e=0.
\]
Therefore
\[
 aI\preceq B\preceq cI,\qquad
 a,c\in\operatorname{spec}(B).
 \tag{BL9}
\]
Compare $B$ with $\operatorname{diag}(J_N,d)$. Their only differences
are at $(q,q+1)$ and $(q+1,q)$, so (BL1) gives equality of rooted
moments through degree $2q+1$. Using (BL3),
\[
 h_0(B^k)_{00}=\int t^k\,d\mu\qquad(0\le k\le2q+1).
 \tag{BL10}
\]
No original Jacobi coefficient after index $q$ is needed.
Again $r(x)=\pi_q(x)/\pi_{q+1}(x)$, and
\[
 \det(tI-B)=(t-d)\pi_{q+1}(t)-\gamma\pi_q(t).
\]
The two endpoint equations therefore also prove the advertised
characteristic-polynomial formula.

\subsection{Nodes, weights, and degenerate cases}

Each nondegenerate border is a real symmetric tridiagonal matrix with
nonzero off-diagonals. If an eigenvector has first coordinate zero,
the eigenvector equations successively force every coordinate to be
zero. Thus each eigenspace has dimension at most one, all eigenvalues
are distinct, and every normalized eigenvector $u_i$ has
$(u_i)_0\ne0$. Spectral diagonalization gives the quadrature rule
\[
 Q(f)=\sum_i w_i f(\theta_i),\qquad
 w_i=h_0|(u_i)_0|^2>0,
 \qquad \sum_i w_i=h_0.
 \tag{BL11}
\]
Equations (BL4), (BL9) put the nodes in $[a,c]$ and include the
prescribed endpoint or endpoints. Equations (BL5), (BL10) prove
the required exactness for every polynomial by linearity.
This proves the Radau and Lobatto border theorems entirely from
finite matrix algebra and the compression identity (BL3).

When a hypothesis of strict endpoint invertibility fails at finite
termination, do not divide by the vanishing determinant. The original
atomic spectral measure already supplies exact quadrature. In the
special Lobatto case $q=0$, the two-endpoint rule remains meaningful
even if a weight vanishes:
\[
 w_a=\int\frac{c-t}{c-a}\,d\mu(t),\qquad
 w_c=\int\frac{t-a}{c-a}\,d\mu(t).
\]
For the zero measure take zero weights. If $a=c$, the one support
point gives the entire integral and no two-distinct-endpoint theorem
is asserted.

\paragraph{Signs retained from the audit.}
For remainder ``integral minus rule'', positive derivatives of order
$2q+1$ make left Radau a lower bound and right Radau an upper bound.
Positive derivatives of order $2q+2$ make Lobatto an upper bound.
The first two directions must not be reversed. The already established
Hermite-interpolation remainder proof needs no change.

\paragraph{Lean reuse and the new lemma.}
The declaration snapshot saved with the earlier audit, at repository
commit \texttt{e5725bb94e5bd8e326d58ad219260b0218e74c9e}, verifies
\code{NLAlib.LanczosDecomp.dotProduct_aeval_mulVec_eq_sum},
\code{NLAlib.LanczosDecomp.cfc_T_apply_zero_zero},
\code{NLAlib.LanczosDecomp.charpoly_T_monic}, and
\code{NLAlib.integral_pow_spectralMeasure}, all without \texttt{sorry}.
They supply the existing moment and spectral-measure interfaces.
The new reusable proof burden is (BL1), followed by the two scalar
Schur-complement calculations. The current local declaration index
inspected for this note does not contain
these Lanczos declarations; the reuse claim is against the specified
verified snapshot, not against that local index.

\section{Christoffel--Uvarov modification by an invertible coefficient system}
\label{focused:modification}

The finite coefficient system is more general than a formula written
in terms of roots. It permits repeated roots, never introduces
negative-index orthogonal polynomials, and makes the normalization
determinant's nonvanishing a one-line positivity argument.

\subsection{A real system that works for every denominator degree}

Let $\mu$ be a finite positive measure on a compact interval $I$ and
$L(f)=\int f\,d\mu$. Let $u,v\in\mathbb R[t]$ be nonzero,
$\ell=\deg u$, $m=\deg v$, and suppose
\[
 v(t)\ne0\ (t\in I),\qquad u(t)/v(t)>0\ (t\in I).
\]
Define $\widehat L(f)=L(uf/v)$ and fix $n\ge0$. Assume
$\widehat L(s^2)>0$ for every nonzero real polynomial $s$ of degree
less than $n$. This condition is vacuous for $n=0$.
Continuity on the compact interval makes every integral below finite.
Common factors of $u$ and $v$ may be cancelled, but cancellation is
not required for the construction.

Set $d=n+\ell$ and $\lambda=\operatorname{lc}(u)\ne0$.
For a polynomial $P$ of degree at most $d$, define the $d$ linear
constraints $\mathcal C(P)$ to be
\[
 \begin{array}{ll}
 [t^i](P\bmod u),&0\le i<\ell,\\[2pt]
 L(t^rP/v),&0\le r<n.
 \end{array}
 \tag{CU1}
\]
Polynomial remainder is defined over the coefficient field, so $u$
need not be monic. When $\ell=0$ there are no remainder constraints.
The desired numerator $P=u\widehat\pi_n$ is characterized by
\[
 \mathcal C(P)=0,\qquad [t^d]P=\lambda.
 \tag{CU2}
\]
Indeed, the first constraints say $u\mid P$; hence $P/u$ is monic
of degree $n$, and the remaining constraints are exactly its
orthogonality for $\widehat L$.

For $0\le j<d$, take the column $\mathcal C(t^j)$ and form the real
$d\times d$ matrix
\[
 M=[\mathcal C(1)\ \mathcal C(t)\ \cdots\ \mathcal C(t^{d-1})],
 \qquad b=\mathcal C(t^d).
\]
The complete computation is the single finite linear system
\[
 M c=-\lambda b,\qquad
 P(t)=\lambda t^d+\sum_{j=0}^{d-1}c_jt^j,
 \qquad \widehat\pi_n=P/u.
 \tag{CU3}
\]

\begin{lemma}[The normalization matrix is nonsingular]
Under the stated positivity hypothesis, $\det M\ne0$.
\end{lemma}
\begin{proof}
Let $S$ have degree less than $d$ and $\mathcal C(S)=0$.
The remainder equations give $S=us$ with $\deg s<n$.
The other equations say $\widehat L(t^rs)=0$ for $0\le r<n$.
Take their linear combination with the coefficients of $s$ to get
$\widehat L(s^2)=0$. Positivity gives $s=0$, hence $S=0$.
Thus the square matrix has zero kernel.
\end{proof}

For $d=0$ the matrix is empty, its determinant is one, and
$P=u$, $\widehat\pi_0=1$. If $n=0$ but $\ell>0$, the remainder
matrix on $1,\ldots,t^{\ell-1}$ is the identity; again $P=u$.
No degree comparison between $m$ and $n$ has entered the proof.

\subsection{The determinant formula, with its exact normalization minor}

Let $D(t)$ be the $(d+1)\times(d+1)$ determinant whose first row is
$(1,t,\ldots,t^d)$ and whose remaining rows are the constraints:
\[
 D(t)=\det\begin{pmatrix}
 1&t&\cdots&t^d\\
 \mathcal C(1)&\mathcal C(t)&\cdots&\mathcal C(t^d)
 \end{pmatrix}.
 \tag{CU4}
\]
Expanding the first row and applying any component of $\mathcal C$
produces a determinant with two equal rows. Thus $\mathcal C(D)=0$.
The highest coefficient is exactly
\[
 \Delta=[t^d]D(t)=(-1)^d\det M\ne0.
\]
Consequently
\[
 \boxed{\quad
 \widehat\pi_n(t)=\frac{\lambda D(t)}{\Delta u(t)}.
 \quad}
 \tag{CU5}
\]
The quotient is polynomial by (CU1), including at zeros of $u$.
This proves the formula and its nonzero normalization, rather than
assuming either. It uses only polynomial division, a finite moment
matrix, and positivity. There is no need to split real polynomials
into complex roots.

\subsection{A unified small determinant, including $m>n$}

If the original monic orthogonal polynomials needed below exist,
the full monomial system can be compressed further. Put
\[
 b_0=\max(n-m,0),\qquad s=\min(n,m),\qquad
 \mathcal V=\operatorname{span}
 \{\pi_{b_0},\pi_{b_0+1},\ldots,\pi_{n+\ell}\}.
 \tag{CU6}
\]
Precisely, assume the displayed polynomials and all preceding ones
are monic and orthogonal, with nonzero squared norms through degree
$b_0-1$. These are the norms used to remove the earlier coefficients.
There are $\ell+s+1$ displayed polynomials. A numerator satisfying
(CU2) belongs to $\mathcal V$: if $n\ge m$, test its orthogonality
with $v r$, $\deg r<n-m$, to get $L(Pr)=0$, which removes the
coefficients of $\pi_0,\ldots,\pi_{n-m-1}$. If $n<m$, the assertion
just says that all polynomials through degree $n+\ell$ are allowed.

On $\mathcal V$, the entire set of orthogonality equations is
equivalent to only $s$ inverse-moment equations:
\[
 P\bmod u=0,\qquad L(t^rP/v)=0\quad(0\le r<s).
 \tag{CU7}
\]
For $n<m$ this is the original system. For $n\ge m$, divide each
test polynomial $h$, $\deg h<n$, as $h=v a+b$ with
$\deg a<n-m$ and $\deg b<m$. Then
$L(hP/v)=L(aP)+L(bP/v)$; its first term vanishes by membership in
$\mathcal V$, and the second by (CU7).

Replace the first row in (CU4) by the basis in (CU6), and the
constraints by the $\ell+s$ rows in (CU7). The determinant now has
size $\ell+s+1$, its degree is at most $n+\ell$, and its leading
coefficient is
\[
 \Delta_{\rm small}=(-1)^{\ell+s}
 \det\!\big[\mathcal C_{\rm small}(\pi_j)\big]_
          {j=b_0}^{n+\ell-1}\ne0.
 \tag{CU8}
\]
To prove nonvanishing, a vector in this minor's kernel produces
$S\in\mathcal V$ of degree less than $n+\ell$ satisfying (CU7).
The equivalence just proved recovers all of (CU1), so the same
positivity argument gives $S=0$.
Formula (CU5) holds with this smaller determinant and normalization.

This is the precise correction to the often displayed index range:
\[
 \begin{array}{c|c|c}
 &\text{polynomial columns}&\text{inverse-moment rows}\\\hline
 m\le n&\pi_{n-m},\ldots,\pi_{n+\ell}&m\\
 m>n&\pi_0,\ldots,\pi_{n+\ell}&n
 \end{array}
\]
For $m>n$, imposing $m$ independent Cauchy-transform equations would
generally overdetermine the problem. One needs the $n$ combinations
in (CU7). If original orthogonal polynomials are unavailable beyond
a finite-support termination index, retain the monomial system
(CU1)--(CU5), whose existence needs only the modified positivity
assumption through degree $n-1$.

For a concrete $m>n$ check, take $\mu=(\delta_{-1}+\delta_1)/2$,
$u=1$, $v=t^2+1$, and $n=1$. The single coefficient equation is
$(1/2)c_0=0$, so $\widehat\pi_1=t$. Its two individual Cauchy
values are both $L(t/(\pm i-t))=-1/2$, rather than zero. The required
combined equation $L(t/(t^2+1))=0$ does hold. Thus the two pole rows
cannot replace the one orthogonality row in this case.

\subsection{Recovering root and Cauchy-transform formulas}

If $u$ has distinct complex roots $u_i$, the remainder rows can be
replaced by $P(u_i)=0$: the Vandermonde matrix makes this an invertible
change of rows. Repeated roots require the corresponding derivative
or Hasse-derivative conditions; retaining polynomial remainders
avoids that extra interface entirely.

If $v$ has distinct complex roots $v_j$, define
$\eta_k(z)=L(\pi_k(t)/(z-t))$. For $0\le r<s\le m$,
\[
 \frac{t^r}{v(t)}=
 \sum_{j=1}^m\frac{v_j^r}{v'(v_j)(t-v_j)},\qquad
 L\!\left(\frac{t^r\pi_k(t)}{v(t)}\right)
 =-\sum_{j=1}^m\frac{v_j^r\eta_k(v_j)}{v'(v_j)}.
 \tag{CU9}
\]
Clear denominators: the polynomial identity has degree less than
$m$ and agrees at the $m$ distinct roots. For $m\le n$, $s=m$,
and the coefficient matrix $(v_j^r/v'(v_j))_{r,j}$ is invertible.
Hence the familiar $m$ separate Cauchy-transform rows are equivalent.
For $m>n$, retain the $n$ displayed combinations, not $m$ separate
rows. Repeated denominator roots are already covered by (CU1);
partial fractions with higher powers are an optional presentation,
not a prerequisite for the theorem.

Complex conjugate roots present no reality problem. The primary
coefficient system (CU3) is real and has a unique solution. Its
root-based versions therefore produce the same real polynomial,
even if their intermediate determinants are complex.

Once these monic polynomials are obtained, no separate modification
algorithm is needed for the recurrence coefficients. If $c_j$ is
the coefficient of $t^{j-1}$ in $\widehat\pi_j$ (with $c_0=0$), then
\[
 \widehat\alpha_j=c_j-c_{j+1},\qquad
 \widehat\beta_j=\widehat h_j/\widehat h_{j-1}\quad(j\ge1),
 \qquad \widehat h_j=L(u\widehat\pi_j^2/v).
\]
The norm ratio requires $\widehat h_{j-1}>0$; an orthonormal
off-diagonal is $\sqrt{\widehat\beta_j}$. This states exactly which
additional moment positivity is used when extending the recurrence.

\subsection{The complex-point nonvanishing correction remains necessary}

For a single Christoffel shift $z$ and monic original polynomials,
the formula is
\[
 \widehat\pi_j(t)=
 \frac{\pi_{j+1}(t)-\pi_{j+1}(z)\pi_j(t)/\pi_j(z)}{t-z},
 \qquad \pi_j(z)\ne0.
 \tag{CU10}
\]
Its numerator is divisible by $t-z$. This does \emph{not} imply
$\widehat\pi_j(z)\ne0$: the numerator may have a double root there.
For the standard Gaussian probability measure,
$\pi_1(t)=t$, $\pi_2(t)=t^2-1$. At $z=i$,
\[
 \widehat\pi_1(t)=\frac{t^2-1-2it}{t-i}=t-i,
 \qquad \widehat\pi_1(i)=0.
 \tag{CU11}
\]
The original Gaussian orthogonal polynomials have real roots, so
$\pi_k(i)\ne0$ for every $k$. Thus even nonvanishing at the shift
for all original polynomials does not repair the false conclusion.
For real $z$ strictly outside the support interval, the modified
functional has constant sign and becomes positive after an overall
sign change; the usual root-location theorem then does give the
extra nonvanishing. That additional real-order hypothesis must be
stated explicitly.

\paragraph{Formalization boundary.}
The coefficient formulation requires no complex-analysis library and
no theorem about sequential modification algorithms. Its reusable
pieces are the polynomial remainder map, finite integrals, the
positive-definite Gram argument, and Cramer's rule. Root splitting
and partial fractions are only corollaries for matching traditional
notation. No complete Christoffel--Uvarov declaration was found in
the inspected local declaration index or in the saved audit snapshot.


\appendix
\section{The retained lower bound for reciprocal approximation}
\label{app:retained-lower}

This appendix makes the exact root-limit conclusion self-contained. It is
the earlier lower-bound route, not an additional simplification of it.
The explicit positive denominator in Section~\ref{sec:explicit-reciprocal}
does not require any of the minimax machinery used here.

Put $f(x)=e^{x/4}$ and use the same weighted inner product and $g_n$ as in
Section~\ref{sec:explicit-reciprocal}. Orthogonality of $f-g_n$ to every
polynomial of degree at most $n$ strengthens the finite Gram calculation to
\begin{equation}
 \int_0^\infty(f-p)^2e^{-x}\,dx
 =\frac2{3^{2n+2}}+\int_0^\infty(p-g_n)^2e^{-x}\,dx
 \ge\frac2{3^{2n+2}}\qquad(\deg p\le n).
 \label{eq:retained-gram-lower}
\end{equation}

Fix $a>0$. Let $q$ minimize
$\max_{[0,a]}e^{-x/2}|f(x)-q(x)|$ over degree-at-most-$n$ polynomials.
Existence follows without an infinite-dimensional compactness theorem:
a bounded sublevel bounds the polynomial at $n+1$ fixed distinct points;
the Vandermonde inverse bounds its coefficients; the objective is continuous
on this closed bounded finite-dimensional set.

The weighted alternation argument gives $n+2$ ordered extrema with alternating
signs. To recall its exact dependency, if there were at most $n+1$ alternating
sign blocks on the compact extremal set, choose one root between each pair
of blocks. A polynomial of degree at most $n$ can then have the sign of the
error everywhere on that set and be nonzero there. A sufficiently small
multiple strictly decreases the weighted error near the extremal set, while
compactness leaves a uniform strict margin on its complement. This contradicts
minimality. The error is nonzero because $f^{(n+1)}>0$ on the interval.

Between those extrema lie $n+1$ distinct zeros
$t_0<\cdots<t_n$ of $f-q$, all in $(0,a)$. Thus $q$ is the Newton interpolant
of $f$ at those points:
\[
 q(x)=\sum_{j=0}^n f[t_0,\ldots,t_j]
                  \prod_{\ell<j}(x-t_\ell).
\]
Every displayed divided difference is positive by repeated Rolle's theorem.
For $x\ge a$, all factors are nonnegative, and the interpolation remainder
is also nonnegative, because $f^{(n+1)}>0$. Consequently
\begin{equation}
 0<q(x)\le f(x)\qquad(x\ge a).
 \label{eq:retained-extrapolation}
\end{equation}

Now set $a=(4n+4)\log3+\log4>2$. Suppose an admissible reciprocal $1/p$
has uniform error
\[
 \epsilon<\frac1{2\sqrt a\,3^{n+1}}
\]
relative to $e^{-x/4}$ on the half-line. Its value at zero and the absence
of poles imply $p>0$ everywhere on the connected half-line. Since
$\epsilon e^{a/4}<1/2$, on $[0,a]$ we have
\[
 p\le2f,\qquad |p-f|\le\epsilon pf\le2\epsilon e^{x/2}.
\]
Minimality of $q$ gives the same weighted error bound for $q$ on this
interval. Combining it with \eqref{eq:retained-extrapolation} and
\eqref{eq:retained-gram-lower} yields
\[
 \frac2{3^{2n+2}}
 \le\int_0^\infty(q-f)^2e^{-x}\,dx
 \le4a\epsilon^2+2e^{-a/2}
 =4a\epsilon^2+\frac1{3^{2n+2}},
\]
contradicting the assumed error. A nonzero constant numerator can be absorbed
into $p$; a zero numerator has error one. Scaling the independent variable
back therefore proves, for every $n\ge0$,
\[
 E_{0,n}(e^{-x};[0,\infty))\ge
 \frac{3^{-n}}{6\sqrt{(4n+4)\log3+\log4}}.
\]
Together with the new explicit upper bound, taking $n$th roots gives $1/3$.
This argument identifies precisely the weighted-minimax component that still
remains if the entire two-sided theorem is to be formalized.

\section{Source and verification notes}
The existing source document is
\code{open_mathematical_items_operator_derivations.tex}; it is not modified
by this separate note. The derivations here replace the specified upper-bound,
degree-transfer, border, and coefficient-system arguments. The retained lower
bound is reproduced only to make its dependency explicit.

For classical context, the reciprocal root limit is recorded in
\href{https://doi.org/10.1137/0719077}{A. Sch\"onhage,
\emph{Rational Approximation to $e^{-x}$ and Related $\mathcal L^2$-Problems}
(1982)}. The polynomial Bernstein limit and its usual entire-function
formulation are discussed in
\href{https://lubinsky.math.gatech.edu/Research\%20papers/BrnstnDec05CA.pdf}
{D. S. Lubinsky, \emph{On the Bernstein Constants of Polynomial Approximation}}.
The finite proofs in this note spell out their own inequalities and
normalizations; no claim of historical novelty is made.

All named NLAlib reuse claims refer to the
\href{https://github.com/tchen-research/NLAlib/blob/e5725bb94e5bd8e326d58ad219260b0218e74c9e/atlas/declarations.json}
{saved public declaration index}. Mathematical review, availability of a
declaration in that snapshot, and a new Lean proof are separate assertions.
Only the first two are claimed here.
\end{document}
